\documentclass[10pt,a4paper]{amsart}
\usepackage[margin=19mm]{geometry}
\usepackage{amsmath,amssymb,mathtools}
\usepackage[scr=rsfs,cal=euler]{mathalfa}
\usepackage{xfrac}
\usepackage[unicode,hidelinks,bookmarks=false]{hyperref}
\newcommand{\ie}{i.e.\ }
\newcommand{\cf}{cf.\ }
\newcommand{\resp}{respectively\ }
\newcommand{\Subsection}{\subsection}
\providecommand{\nobreakdash}{\nobreak\hbox{-}}
\setlength{\emergencystretch}{2em}
\multlinegap0pt
\usepackage{bm}
\usepackage{bbm}
\usepackage{dsfont}
\usepackage{xspace}
\usepackage{cancel}
\usepackage{mathtools}
\usepackage{cleveref}
\usepackage{amsthm}
\makeatletter
\@ifundefined{theorem}{\newtheorem{theorem}{Theorem}}{}
\@ifundefined{proposition}{\newtheorem{proposition}[theorem]{Proposition}}{}
\makeatother
\newcommand{\mC}{\mathcal{C}}
\newcommand{\mR}{\mathcal{R}}
\newcommand{\mU}{\mathcal{U}}
\title[Linear Code Conversion in the Merge Regime: General Bounds a]{Linear Code Conversion in the Merge Regime: General Bounds and Reed-Muller Constructions}
\author{Anina Gruica, Benjamin Jany, Stanislav Kruglik}
\date{Source: arXiv:2606.27179v2 (CC BY 4.0)}
\begin{document}
\maketitle

\noindent\emph{Source and licence.} Linear Code Conversion in the Merge Regime: General Bounds and Reed-Muller Constructions. Version arXiv:2606.27179v2 (CC BY 4.0), distributed under the Creative Commons Attribution 4.0 International licence, as recorded in the arXiv licence metadata for this exact version and shown on the abstract page https://arxiv.org/abs/2606.27179v2. Full rights notice: licenses/arxiv-2606.27179-CC-BY-4.0.txt.\par\smallskip

\noindent\textit{Editorial scope.} Counted unit: the Theorem whose source label is \texttt{thm:readlowergen} in the article (source0.tex lines 716--719, source label \texttt{thm:readlowergen}), delivered together with the article's own complete original proof of it (source0.tex lines 722--765, one \texttt{proof} environment of 44 lines). No mathematics is written, repaired or re-derived: statement and proof are the article's own text line for line. Required context retained uncounted: prop:readlower (lines 697--701). Every definition, display and label of those lines is retained unchanged. Omitted from this package and not redistributed: 1--696, 702--715, 720--721, 766--1339. Nothing inside a retained excerpt is omitted or altered: every retained body is delivered exactly as the article's lines are written, so no omission receipt of any kind accompanies this package. Citations: the retained excerpts carry no citation command at all, so no bibliography entry of the article is transcribed and no reference is redistributed. Reference adaptation: none was needed and none was made; every reference inside the delivered unit resolves inside it, so no unresolved reference or citation is produced. Printed slips kept literal: no printed word, symbol, hypothesis or number of the counted statement or of its proof was altered. Layout and class: the article's own class is \texttt{IEEEtran} with packages including setspace, amsmath, amssymb, mathrsfs, amsthm, multirow, color, array, booktabs, float, url, comment, enumerate, graphicx; the wrapper is the lane's pinned \texttt{amsart} editorial wrapper at 10pt on A4, which restates the theorem-like environments the retained text needs and adds the article's own macro definitions that the retained text needs (\texttt{\textbackslash mC}, \texttt{\textbackslash mR}, \texttt{\textbackslash mU}), with extra wrapper packages (bm, bbm, dsfont, xspace, cancel, mathtools, cleveref). The delivered numbering is the unit's own. Assets: no retained span contains a figure environment, a graphics-inclusion command, a PDF-page inclusion, a table or a TikZ picture; no image file is copied into this package, the package declares no asset dependency and no image file is redistributed with it. Licence: version arXiv:2606.27179v2 of this article is distributed under the Creative Commons Attribution 4.0 International licence, as recorded in the arXiv licence metadata for this exact version and shown on the abstract page https://arxiv.org/abs/2606.27179v2. A full rights notice is delivered at licenses/arxiv-2606.27179-CC-BY-4.0.txt. Characters: every retained excerpt is pure ASCII, so the delivered engine, XeLaTeX with its default Latin Modern text font, reports no missing character.
\begin{proposition}\label{prop:readlower}
Let $(\mC_{\textbf{I}},\mC_F,\sigma)$ be a convertible code. Fix $i\in [\lambda]$, and let $\mR_i$ and $\mU_i$ be the sets of read and unchanged symbols from $\mC_{I_i}$, respectively. Then
\begin{align*} |\mR_i|\ge \max_{1\le j\le k_{I_i}}\left\{ k_{I_i}-j+1-\bigl(|\mU_i|-d_j(\mC_F)+1\bigr)^+ \right\}, \end{align*}
where $(a)^+:=\max\{a,0\}$.
\end{proposition}

\begin{theorem}\label{thm:readlowergen}

Let $(\mC_{\mathbf I},\mC_F,\sigma)$ be a convertible code. Fix $i\in[\lambda]$, and let $\mR_i$ and $\mU_i$ be the sets of read and unchanged symbols from $\mC_{I_i}$, respectively. If there exists an index $j\in[k_{I_i}]$ such that $d_j(\mC_F)<|\mU_i|+1$, let $j$ be the largest such index. Then \[ |\mR_i| \geq \begin{cases} k_{I_i}-j, & \text{if such an index } j \text{ exists},\\ k_{I_i}, & \text{otherwise}. \end{cases} \]
\end{theorem}

\begin{proof}
If no index $j\in[k_{I_i}]$ satisfies $d_j(\mC_F)<|\mU_i|+1$, then
$d_\ell(\mC_F)\ge |\mU_i|+1$ for all $\ell\in[k_{I_i}]$. Hence $(|\mU_i|-d_\ell(\mC_F)+1)^+=0$ for all $\ell$, and the maximum in Proposition~\ref{prop:readlower} is achieved at $\ell=1$, with value $k_{I_i}$. 

Now assume that such an index exists, and let $j\in[k_{I_i}]$ be the largest index such that
$d_j(\mC_F)<|\mU_i|+1$. If $j=k_{I_i}$, then the desired bound is
$|\mR_i|\ge k_{I_i}-j=0$, which is trivial. Thus, in the rest of the proof,
we may assume that $j<k_{I_i}$.

We compare the terms in the maximum from \cref{prop:readlower}.
For $\ell\le j$, we have
$|\mU_i|-d_\ell(\mC_F)+1\ge |\mU_i|-d_j(\mC_F)+1>0$. Moreover, by strict
monotonicity of the generalized Hamming weights,
$d_j(\mC_F)-d_\ell(\mC_F)\ge j-\ell$. Hence
\begin{align*}
&k_{I_i}-j+1-\bigl(|\mU_i|-d_j(\mC_F)+1\bigr)^+\\
&=
k_{I_i}-j+1-|\mU_i|+d_j(\mC_F)-1 \\
&\ge
k_{I_i}-j+1-|\mU_i|+j-\ell+d_\ell(\mC_F)-1 \\
&=
k_{I_i}-\ell+1-\bigl(|\mU_i|-d_\ell(\mC_F)+1\bigr)^+ .
\end{align*}
Thus, among all $\ell\le j$, the largest possible value is bounded above by the term with $\ell=j$.

Since $j$ is maximal and $j<k_{I_i}$, we have
$d_\ell(\mC_F)\ge |\mU_i|+1$ for all $\ell>j$. Therefore
$(|\mU_i|-d_\ell(\mC_F)+1)^+=0$ for all $\ell>j$, and hence $k_{I_i}-\ell+1-\bigl(|\mU_i|-d_\ell(\mC_F)+1\bigr)^+
=
k_{I_i}-\ell+1
\le
k_{I_i}-j$ for all $\ell>j$. In particular, for $\ell=j+1$, we obtain $k_{I_i}-(j+1)+1-\bigl(|\mU_i|-d_{j+1}(\mC_F)+1\bigr)^+
=
k_{I_i}-j$.

It remains to compare this value with the maximum over $\ell\le j$. Since
$d_j(\mC_F)<|\mU_i|+1$ and all quantities are integers, we have
$(|\mU_i|-d_j(\mC_F)+1)^+\ge 1$. Therefore $k_{I_i}-j+1-\bigl(|\mU_i|-d_j(\mC_F)+1\bigr)^+
\le
k_{I_i}-j$.

Combining the inequalities above, the term with $\ell=j+1$ is at least as large as all other terms in the maximum from \cref{prop:readlower}, and its value is $k_{I_i}-j$. Hence
$|\mR_i|\ge k_{I_i}-j$, as claimed.
    \end{proof}

\end{document}
